%%
%% UTF-8       UTF-8      UTF-8       UTF-8       UTF-8       UTF-8
%% UTF-8       UTF-8      UTF-8       UTF-8       UTF-8       UTF-8
%%
%% XeLaTeX  XeLaTeX  XeLaTeX  XeLaTeX  XeLaTeX  XeLaTeX  XeLaTeX  XeLaTeX  

% !TeX program = xelatex
\lccode`\-=`\-
%\defaulthyphenchar=127
\documentclass[12pt,a4paper,oneside]{extarticle}
%\usepackage{xunicode} % https://ctan.org/pkg/xunicode
%\usepackage{xltxtra} % https://ctan.org/pkg/xltxtra
\usepackage{fontspec} % https://ctan.org/pkg/fontspec
\usepackage{polyglossia} % https://ctan.org/pkg/polyglossia
\setdefaultlanguage{portuges}
%\setdefaultlanguage{portuguese}
%\setdefaultlanguage[spelling=old]{portuguese}
\usepackage{csquotes}

%\usepackage{ebgaramond}%[RawFeature={+onum,+pnum,+ss06}]
\usepackage[cmintegrals,cmbraces]{newtxmath}
%\usepackage{ebgaramond-maths}

\usepackage[no-math]{fontspec}
%\setmainfont{EBGaramond}[
%  Extension=.otf,
%  UprightFont=*-Regular,
%  ItalicFont=*-Italic,
%  BoldFont=*-Bold,
%  BoldItalicFont=*-BoldItalic,
%  Numbers={Proportional,OldStyle},
%  Ligatures=TeX
%]

%\setmainfont{TeX Gyre Pagella}[
%  Extension = .otf,
%  UprightFont = *-Regular,
%  ItalicFont = *-Italic,
%  BoldFont = *-Bold,
%  BoldItalicFont = *-BoldItalic,
%  Ligatures = TeX
%]



\clubpenalty=100
\linepenalty=100
\widowpenalty=100

%estética
% \usepackage[hyphens]{url}
\usepackage{xurl}
\urlstyle{same} % URLs na mesma fonte do texto (Times New Roman)
\AtBeginDocument{%
  \setlength{\Urlmuskip}{0mu plus 1mu} % Ajusta espaçamento entre caracteres em URLs
  \setlength{\emergencystretch}{3em} % Permite maior flexibilidade na justificação
  \def\UrlBreakPenalty{100} % Penalidade baixa para facilitar quebras
  \def\UrlBigBreakPenalty{50} % Penalidade para quebras em caracteres como /
  \sloppy % Relaxa a justificação para evitar ultrapassagem
}

\usepackage[hang,small,bf]{caption}
\usepackage[a4paper,left=2.5cm,right=2.5cm,top=2.5cm,bottom=2.5cm,includehead,includefoot]{geometry}
\linespread{1.3}

\usepackage[
  xetex,                      % Driver para XeLaTeX
  unicode=true,               % Suporte a caracteres Unicode (ex.: ç, ã)
  bookmarks=true,             % Cria bookmarks no PDF
  bookmarksopen=true,         % Bookmarks abertos por padrão
  bookmarksnumbered=true,     % Numera os bookmarks (ex.: 1, 1.1, 1.2)
  linktoc=all,                % Links clicáveis em sumário, figuras e citações
  breaklinks=true,            % Permite quebra de links longos
  colorlinks=true,            % Links com cor em vez de caixas
  citecolor=black,            % Cor das citações (preto)
  filecolor=black,            % Cor dos links de arquivos (preto)
  linkcolor=black,            % Cor dos links internos (preto)
  urlcolor=blue,              % URLs em azul
  pdfpagelabels=true,         % Adiciona rótulos às páginas do PDF
  plainpages=false,           % Diferencia numeração de páginas
  pdftitle={Desafios 12 ano},
  pdfauthor={Desafios 12 ano}, % Nome do autor
  pdfsubject={Desafios 12 ano},
  pdfkeywords={Desafios 12 ano},
  pdfcreator={XeLaTeX com hyperref},
  pdfproducer={XeLaTeX}
]{hyperref}

%% Definir um novo estilo de url
\makeatletter
\def\url@luisstyle{%
  \@ifundefined{selectfont}{\def\UrlFont{\sf}}{\def\UrlFont{\small\ttfamily}}}
\makeatother
\urlstyle{luis}

\setlength{\captionmargin}{1.5in}
\usepackage{fancyhdr}

\renewcommand{\thefootnote}{\alph{footnote}} %modificar as letras de notas de rodapé%
\makeatletter
\renewcommand\@makefnmark{\@textsuperscript{\normalfont(\@thefnmark)}}
\makeatother
\renewcommand{\fboxrule}{2pt} %para alterar a expessura das caixas
%\renewcommand{\fboxsep}{3ex} %para alterar a separação nas caixas

%para alterar a fonte dos items http://cnlart.web.cern.ch/cnlart/214/node68.html
%\renewcommand{\labelenumi}{\textbf{\arabic{enumi}.}}
\renewcommand{\labelenumi}{\textbf{\alph{enumi}.}}
%\renewcommand{\labelenumii}{\textbf{\alph{enumii}.}}
%\renewcommand{\labelenumii}{\textbf{(\alph{enumii})}}
\renewcommand{\labelenumii}{\textbf{(\arabic{enumii})}}

\long\def\symbolfootnote[#1]#2{\begingroup\def\thefootnote{\fnsymbol{footnote}}\footnote[#1]{#2}\endgroup}

%%%%%%%%%%%%%%%%%%%%%%%%%%%%
%%%%%%%%%%%%%%%%%%%%%%%%%%%%
%%%%%%%%%%%%%%%%%%%%%%%%%%%%
\usepackage{setspace}
%\singlespacing
%\onehalfspacing
%\doublespacing

\usepackage{comment}
\usepackage{siunitx}
%%%%%%%%%%%%%%%%%%%%%%%%%%%%
\usepackage[x11names]{xcolor} 
%%%%%%%%%%%%%%%%%%%%%%%%%%%%
\usepackage{relsize}
%\newcommand\CC{C\nolinebreak\hspace{-.05em}\raisebox{.4ex}{\relsize{-3}{\textbf{+}}}\nolinebreak\hspace{-.10em}\raisebox{.4ex}{\relsize{-3}{\textbf{+}}}}
\newcommand\CC{C\nolinebreak[4]\hspace{-.05em}\raisebox{.4ex}{\relsize{-3}{\textbf{++}}}}

%\usepackage{showframe}
\usepackage{pgf,tikz}       % - Elementos gráficos
\usepackage{tikzpagenodes}
\usetikzlibrary{backgrounds,calc,positioning}
\usetikzlibrary{positioning, shapes}

\pgfdeclarelayer{background}
\pgfdeclarelayer{foreground}
\pgfsetlayers{background,main,foreground}

\usepackage[
    backend=biber,
    style=numeric-comp,
    language=auto,
    natbib=false,
    block=space,
    isbn=false,
    url=true, 
    doi=true,
    mcite=true,
    eprint=false,
    sorting=none,
    maxnames=6
]{biblatex} % https://ctan.org/pkg/biblatex
\addbibresource{bibliografia.bib}
 % não é para alterar

\begin{document}
%espaço negativo para o cabeçalho
\vspace*{-2cm}
\begin{tikzpicture}[remember picture, overlay]
  \draw[line width = 0.5pt, color=black] ($(current page.north west) + (2.5cm,-2cm)$) -- ($(current page.north east) + (-2.5cm,-2cm)$);
\end{tikzpicture}
\vspace*{2mm}
\begin{minipage}[c][2.5cm][t]{.25\textwidth}
    \includegraphics[width=\linewidth]{logos/aeB06.pdf}
\end{minipage}\hfill
\begin{minipage}[c][2.5cm][t]{.7\textwidth}
	\begin{center}\Large\textbf{Agrupamento de Escolas de Benfica}\end{center}
	\vspace{-8.5mm}
	\begin{center}\large\textbf{Escola Secundária José Gomes Ferreira}\end{center}
	\vspace{-8.5mm}
	\begin{center}\large{Grupo Disciplinar 510}\end{center}
\end{minipage}
\begin{tikzpicture}[remember picture, overlay]
\draw[line width = 0.5pt, color=black] ($(current page.north west) + (2.5cm,-4.8cm)$) -- ($(current page.north east) + (-2.5cm,-4.8cm)$);
\end{tikzpicture}
\vspace{10mm}
\begin{center}
\Large{\textbf{Análise e Implementação Computacional da Geometria dos Favos de Mel}}\par 
\large{\textbf{Adrian Dias}}\par 
\large{\textbf{Nº 1, T 12º 3ª }}\par 
\Large{-- Disciplina de Física --}\par
\textbf{2025 - 2026}\\
\end{center}

\begin{tikzpicture}[remember picture, overlay]
  \draw[line width = 0.5pt, color=black] ($(current page.south west) + (2.5cm,3cm)$) -- ($(current page.south east) + (-2.5cm,3cm)$);
\end{tikzpicture}
 % pode alterar o título, o nome o número etc.

\begin{abstract}
Este relatório apresenta o desenvolvimento de um modelo matemático e a respectiva implementação em \texttt{Python}/\texttt{turtle}, com vista à representação de uma curva ou forma encontrada na natureza. São descritos os objectivos, a metodologia seguida e os principais resultados obtidos. 
\end{abstract}

\section{Introdução}
Neste capítulo apresenta-se o enquadramento teórico do trabalho, incluindo o modelo matemático de base e o modelo específico implementado no código.  
De acordo com a proposta, seleccionou-se uma forma presente na natureza, a estrutura hexagonal dos favos de mel, a qual serve de inspiração para a construção do script em \texttt{Python}/\texttt{turtle}.

\subsection{Conjectura do Favo de Mel e Otimização (1.1)}
A estrutura dos favos de mel obedece à \textbf{Conjectura do Favo de Mel} (Honeycomb Conjecture), demonstrada matematicamente por Thomas Hales em 1999. Esta conjectura estabelece que, entre todos os possíveis padrões de divisão do plano, a configuração hexagonal regular minimiza o perímetro total para uma área fixa. A razão ótima perímetro-área é dada por:

\begin{equation}
\min \frac{P}{A} = \sqrt[4]{12} \approx 1,8612
\label{eq:honeycomb_conjecture}
\end{equation}

onde $P$ representa o perímetro e $A$ a área. No caso do hexágono regular com lado $s$, temos:

\begin{equation}
\frac{P}{A} = \frac{6s}{\frac{3\sqrt{3}}{2}s^2} = \frac{4}{\sqrt{3}s}
\label{eq:hexagon_ratio}
\end{equation}

\subsection{Limitações das Tesselações Regulares (1.2)}
A geometria euclidiana impõe restrições rigorosas às tesselações regulares do plano. Conforme demonstrado matematicamente, apenas três polígonos regulares preenchem completamente o espaço sem sobreposições ou espaços vazios:

\begin{itemize}
    \item \textbf{Triângulos equiláteros} (notação de Schläfli: $\{3,6\}$)
    \item \textbf{Quadrados} (notação de Schläfli: $\{4,4\}$) 
    \item \textbf{Hexágonos regulares} (notação de Schläfli: $\{6,3\}$)
\end{itemize}

A condição matemática para tesselação regular é dada por:

\begin{equation}
\frac{1}{m} + \frac{1}{n} = \frac{1}{2}
\label{eq:tessellation_condition}
\end{equation}

onde $m$ e $n$ são inteiros que satisfazem a equação para os casos $\{3,6\}$, $\{4,4\}$ e $\{6,3\}$.

\subsection{Impossibilidade de Refinamento Hierárquico (1.3)}
Os hexágonos regulares apresentam limitações matemáticas inerentes no que concerne a sistemas de múltiplas escalas. A impossibilidade de subdivisão hierárquica perfeita deve-se a três restrições fundamentais:

\begin{enumerate}
    \item \textbf{Conservação de área}: $A_{\text{pai}} = \sum A_{\text{filhos}}$
    \item \textbf{Hierarquia simples}: Apenas um pai por célula
    \item \textbf{Cobertura perfeita}: Sem buracos ou sobreposições
\end{enumerate}

Estas restrições matemáticas explicam por que as abelhas constroem favos de tamanho uniforme em vez de sistemas com múltiplas escalas.

\subsection{Modelo de Grade Hexagonal Implementado (1.4)}
No script \texttt{Python/turtle}, a estrutura hexagonal é implementada através de um sistema de coordenadas baseado em vetores de rede hexagonal. As posições dos hexágonos são calculadas por:

\begin{equation}
\vec{r}_{ij} = i \cdot \vec{a}_1 + j \cdot \vec{a}_2
\label{eq:hexagonal_grid}
\end{equation}

onde os vetores base da rede hexagonal são definidos como:

\begin{align}
\vec{a}_1 &= \left(1.5s, 0\right) \\
\vec{a}_2 &= \left(0.75s, \frac{\sqrt{3}}{2}s\right)
\label{eq:base_vectors}
\end{align}

sendo $s$ o comprimento do lado do hexágono.

\subsection{Parâmetros e Especificações do Modelo (1.5)}
Para a implementação computacional, foram estabelecidos os seguintes parâmetros matemáticos:

\begin{itemize}
    \item \textbf{Tamanho do hexágono}: $s = 25$ unidades turtle
    \item \textbf{Número de camadas}: $n = 7$ camadas concêntricas
    \item \textbf{Total de hexágonos}: $N = 1 + \sum_{k=1}^{7} 6k = 127$ favos
    \item \textbf{Ângulo de rotação}: $\theta = 60^\circ$ entre lados consecutivos
    \item \textbf{Área individual}: $A_{\text{hex}} = \frac{3\sqrt{3}}{2}s^2$
\end{itemize}

\subsection{Algoritmo de Construção Progressiva (1.6)}
O modelo implementado simula a construção natural através de um algoritmo que segue padrão de expansão radial, iniciando em pontos centrais e expandindo concentricamente. A progressão obedece à sequência:

\begin{equation}
N_k = 6k \quad \text{para} \quad k = 1, 2, \ldots, n
\label{eq:progressive_growth}
\end{equation}

onde $N_k$ representa o número de hexágonos na camada $k$-ésima, garantindo crescimento simétrico e geometricamente consistente.











\section{Parte Experimental}
A componente experimental deste trabalho corresponde à elaboração e explicação do código\cite{carvalho2021praticas}.  
Tal como num procedimento laboratorial, importa detalhar a lógica implementada\cite{martins2015programacao}, os algoritmos utilizados e as opções tomadas em cada etapa, de forma a permitir a replicação do processo.  

%Isto está a dar imensos problemas!!! :( 


\section{Discussão dos Resultados}
Apresentam-se e analisam-se, nesta secção, as imagens geradas automaticamente pelo código.  
Não foram utilizadas capturas de ecrã, mas sim exportações directas produzidas pelo programa.  
Discutem-se as semelhanças e diferenças entre os resultados e a imagem de referência, identificando as causas dos desvios e avaliando a qualidade da aproximação obtida.  

\section{Conclusões}
As conclusões são redigidas a partir da análise dos resultados.  
Evitam-se afirmações superficiais ou subjectivas; privilegiam-se observações fundamentadas, como, por exemplo:  
\begin{itemize}
    \item O modelo reproduz com fidelidade parcial a forma natural seleccionada.  
    \item As limitações decorrem de aproximações matemáticas ou restrições do ambiente de programação.  
    \item Futuras melhorias poderão incluir optimizações algorítmicas ou refinamentos gráficos.  
\end{itemize}

\section{Bibliografia}
\printbibliography


\label{fim}\end{document}
